\chapter{Compilers: From Human Text to Optimizing Silicon}
\label{ch41:compilers}

In Chapter~37, we saw how
CPU pipelines, cache hierarchies, and DRAM bus cycles dictate real execution
speed. But software engineers rarely write raw binary machine instructions or
manually coordinate pipeline registers. Programs are written with
variable names, arithmetic expressions, and control structures such as loops.
The processor that runs them has none of these. It executes binary instructions
on billions of transistors switching several billion times a second, and it
knows only how to fetch a word from memory, decode its bit pattern into control
signals, and write an arithmetic result into a register.

How does source text become a sequence of such machine instructions?

The program that performs this translation is a compiler. A modern optimizing
compiler is a five-stage pipeline of data structures and graph algorithms, and
earlier chapters of this book have already covered the main algorithm behind
each stage:
\begin{itemize}
  \item \textbf{Stage 1: Lexing} scans raw character streams into token
  sequences using the \emph{Finite Automata} of Chapter~31.
  \item \textbf{Stage 2: Parsing} turns flat tokens into an unambiguous
  hierarchical \emph{Tree} (Chapters~10
  and~12) through recursive descent and precedence climbing.
  \item \textbf{Stage 3: Semantic Analysis} checks types and resolves variable
  scopes using an environment stack of \emph{Hash Tables} (Chapter~11).
  \item \textbf{Stage 4: Intermediate Representation \& Optimization} lowers
  trees into three-address code, builds a \emph{Control-Flow Graph} of basic
  blocks, renames variables into \emph{Static Single Assignment (SSA)} form to
  turn data dependencies into an explicit \emph{Directed Acyclic Graph}
  (Chapter~22), eliminates duplicate computations via
  hash tables, and hoists loop invariants.
  \item \textbf{Stage 5: Code Generation \& Register Allocation} maps an
  infinite supply of virtual variables down to a handful of physical CPU
  registers by solving the \emph{Graph Coloring Problem}
  (Chapter~33) via Kempe's degree-peeling heuristic,
  and selects optimal instruction sequences using \emph{Dynamic Programming}
  (Chapter~15).
\end{itemize}

We will follow a single four-line loop snippet throughout this chapter (with \texttt{i} initialized to \texttt{0} on entry):
\begin{lstlisting}[language=C]
sum = 0
while (i < n) {
    sum = sum + a[i] * 2
    i = i + 1
}
\end{lstlisting}
Starting from the raw ASCII bytes on disk, we will trace these four lines by hand:
lexing them into tokens, parsing an abstract syntax tree, lowering types across
lexical scopes, converting control flow into SSA $\phi$-functions, applying four
classical loop optimizations, and finally coloring the liveness interference graph
down to a six-instruction register-only x86-64 loop.

\chapterline

\section{History and Overview}
\label{ch41:setup}

\begin{historybox}
In 1952, Grace Murray Hopper was working at Remington Rand on the UNIVAC I.
Programs were written directly in octal machine code, with programmers looking
up operation codes and memory addresses in thick binders. Hopper realized that
programmers spent days re-deriving routine operations: evaluating sines, taking
square roots, or copying memory blocks. She built the A-0 system, which took a
sequence of call descriptors, searched a magnetic tape library of subroutines,
adjusted absolute addresses, and linked them into an executable routine. Hopper
coined the word ``compiler'' because her program compiled subroutines from a
library. When she presented it, contemporaries scoffed, insisting that machines
were calculating engines designed for arithmetic, not writing programs.

Between 1954 and 1957, John Backus led a small team at IBM to develop FORTRAN
(Formula Translating System) for the IBM 704. In an era where mainframes cost
millions of dollars and CPU time was billed by the minute, hand-written assembly
was considered indispensable. Skeptics insisted that machine-generated
code would run at least five times slower than human assembly, making high-level
languages an unaffordable extravagance. Backus understood that unless FORTRAN's
output rivaled handcrafted assembly in execution speed, the project would be
abandoned. This performance mandate forced Backus's team to invent compiler
optimization from the ground up: they created the first algorithms for loop
nest analysis, execution frequency weighting, and index register allocation.
When FORTRAN was delivered, its emitted code ran with such efficiency that
programmers rapidly left manual assembly behind.

The theoretical foundations took shape in the late 1950s. In 1956, Noam Chomsky
formulated his hierarchy of formal grammars, showing that regular languages
(recognized by finite automata) are strictly less expressive than context-free
languages (which require a stack). When John Backus designed the grammar of
ALGOL 58, he introduced a formal metalanguage, which Peter Naur refined for the
landmark ALGOL 60 report into Backus--Naur Form (BNF). BNF replaced ambiguous
English prose specifications with mathematically rigorous substitution rules,
making parsing a systematic, grammar-driven task.

During the 1960s and 1970s at IBM Research, Frances Allen and John Cocke created
the architectural framework of modern optimizing compilers. Allen formalized the
control-flow graph, introduced interval analysis and dominator trees, and built
the mathematical theory of dataflow equations. Her 1970 paper ``Control Flow
Analysis'' and her joint work with Cocke established the classical repertoire of
loop optimizations---including loop-invariant code motion and induction variable
strength reduction---that every modern compiler performs. Allen was later
awarded the ACM A.M. Turing Award for this work.

Yet for three decades, register allocation remained an ad-hoc heuristic. A CPU
possesses only a handful of hardware registers, while programs manipulate scores
of variables. In 1981, Gregory Chaitin at IBM showed
that assigning $K$ physical registers to variables without interference is
mathematically isomorphic to the graph vertex coloring problem investigated by
Alfred Kempe a century earlier. While graph coloring is NP-complete, Chaitin
showed that a degree-peeling heuristic colors real-world program graphs near-optimally
in linear time, giving compilers their first principled register allocator.

In 1989, Ron Cytron, Jeanne Ferrante, Barry Rosen, Mark Wegman, and Kenneth
Zadeck introduced Static Single Assignment (SSA) form. By placing $\phi$-nodes at
control-flow joins, SSA ensured that every variable is assigned at exactly one
syntactic site. This decoupled values from memory cells, turning tangled dataflow
into clean directed acyclic graphs and accelerating optimization passes from
quadratic to linear time. In 2004, Chris Lattner and Vikram Adve at UIUC unveiled
LLVM, establishing a language-independent, typed SSA intermediate representation.
Today, LLVM, GCC, and just-in-time engines such as Google's V8 and PyTorch's
\texttt{torch.compile} build directly upon this foundation.
\end{historybox}

\begin{focusbox}
A compiler lowers representations through five distinct transformations:
\[
  \text{Characters} \xrightarrow[\text{Lex}]{\text{FA}}
  \text{Tokens} \xrightarrow[\text{Parse}]{\text{Grammar}}
  \text{AST} \xrightarrow[\text{Check}]{\text{Types}}
  \text{SSA IR} \xrightarrow[\text{Optimize}]{\text{Graphs}}
  \text{SSA IR} \xrightarrow[\text{Backend}]{\text{Coloring}}
  \text{Assembly}
\]
Each stage drops surface syntax and makes more of the program's structure explicit.
\end{focusbox}

Figure~\ref{ch41:fig-compiler-pipeline} illustrates this three-tiered architecture. By interposing a shared intermediate representation between languages and targets, modern compilers solve the classic $N \times M$ proliferation problem: $N$ language frontends compile into one shared IR, which undergoes machine-independent optimization before $M$ backends emit native machine code.

\begin{figure}[H]
\centering
\begin{tikzpicture}[
  font=\small,
  tierbox/.style={draw=gray!60, rounded corners=5pt, inner sep=6pt, line width=0.8pt},
  stagebox/.style={draw=gray!70, fill=white, rounded corners=3pt, minimum width=3.6cm, minimum height=0.7cm, align=center, font=\sffamily\scriptsize},
  arrow/.style={-Stealth, thick, gray!70},
  tierarrow/.style={-Stealth, line width=1.5pt, gray!80}
]

% Tier 1: Frontend
\node[tierbox, fill=sourceblue!6, minimum width=3.8cm, minimum height=5.2cm] (fe_box) at (-4.9, 0) {};
\node[anchor=north, font=\sffamily\bfseries\footnotesize, text=sourceblue!90!black] at (-4.9, 2.45) {Frontend (Source)};

\node[stagebox, draw=sourceblue!80] (fe_in) at (-4.9, 1.6) {\textbf{Source Characters}\\ \texttt{\scriptsize sum = 0; while...}};
\node[stagebox, draw=sourceblue!80] (fe_mid) at (-4.9, 0.2) {\textbf{Token Stream}\\ \texttt{\scriptsize IDENT, OP\_ASSIGN, ...}};
\node[stagebox, draw=sourceblue!80] (fe_out) at (-4.9, -1.2) {\textbf{Abstract Syntax Tree}\\ \texttt{\scriptsize While(cond, Block)}};

\draw[arrow] (fe_in) -- node[midway, right=1pt, font=\sffamily\tiny, text=gray!80] {Lexer (DFA)} (fe_mid);
\draw[arrow] (fe_mid) -- node[midway, right=1pt, font=\sffamily\tiny, text=gray!80] {Parser (Grammar)} (fe_out);

\node[font=\sffamily\tiny\itshape, text=sourceblue!80!black] at (-4.9, -2.15) {$N$ Frontends (C, Rust, Swift)};

% Tier 2: Middle-End
\node[tierbox, fill=targetgreen!6, minimum width=3.8cm, minimum height=5.2cm] (me_box) at (0, 0) {};
\node[anchor=north, font=\sffamily\bfseries\footnotesize, text=targetgreen!90!black] at (0, 2.45) {Middle-End (Optimizer)};

\node[stagebox, draw=targetgreen!80] (me_in) at (0, 1.6) {\textbf{Decorated AST / 3AC}\\ \texttt{\scriptsize load(base\_a + i * 8)}};
\node[stagebox, draw=targetgreen!80] (me_mid) at (0, 0.2) {\textbf{SSA Form with $\phi$-Nodes}\\ \texttt{\scriptsize sum\_1 = $\phi$(sum\_0, sum\_2)}};
\node[stagebox, draw=targetgreen!80] (me_out) at (0, -1.2) {\textbf{Optimized SSA CFG}\\ \texttt{\scriptsize limit\_0, ptr + 8, $\ll$1}};

\draw[arrow] (me_in) -- node[midway, right=1pt, font=\sffamily\tiny, text=gray!80] {SSA Placement} (me_mid);
\draw[arrow] (me_mid) -- node[midway, right=1pt, font=\sffamily\tiny, text=gray!80] {LICM, DCE, Folding} (me_out);

\node[font=\sffamily\tiny\itshape, text=targetgreen!80!black] at (0, -2.15) {Universal Optimizer};

% Tier 3: Backend
\node[tierbox, fill=quotientorange!6, minimum width=3.8cm, minimum height=5.2cm] (be_box) at (4.9, 0) {};
\node[anchor=north, font=\sffamily\bfseries\footnotesize, text=quotientorange!90!black] at (4.9, 2.45) {Backend (Machine)};

\node[stagebox, draw=quotientorange!80] (be_in) at (4.9, 1.6) {\textbf{Instruction DAG}\\ \texttt{\scriptsize cmp, jb, add, shl}};
\node[stagebox, draw=quotientorange!80] (be_mid) at (4.9, 0.2) {\textbf{Interference Graph}\\ \texttt{\scriptsize $K$-colorable RIG}};
\node[stagebox, draw=quotientorange!80] (be_out) at (4.9, -1.2) {\textbf{Emitted Assembly}\\ \texttt{\scriptsize x86-64 / ARM64}};

\draw[arrow] (be_in) -- node[midway, right=1pt, font=\sffamily\tiny, text=gray!80] {Liveness Analysis} (be_mid);
\draw[arrow] (be_mid) -- node[midway, right=1pt, font=\sffamily\tiny, text=gray!80] {Chaitin Coloring} (be_out);

\node[font=\sffamily\tiny\itshape, text=quotientorange!80!black] at (4.9, -2.15) {$M$ Backends (x86, ARM)};

% Connecting arrows between tiers
\draw[tierarrow, sourceblue!80] (-2.95, 0.2) -- node[midway, above=1pt, font=\sffamily\tiny\bfseries, text=sourceblue!90!black, align=center] {Type\\Lowering} (-1.95, 0.2);
\draw[tierarrow, targetgreen!80] (1.95, 0.2) -- node[midway, above=1pt, font=\sffamily\tiny\bfseries, text=targetgreen!90!black, align=center] {Target\\Codegen} (2.95, 0.2);

\end{tikzpicture}
\caption{The three-tiered pipeline of modern optimizing compilers. The source-specific Frontend parses characters into a typed AST; the target-independent Middle-End optimizes code in SSA form via loop-invariant motion, strength reduction, and dead code elimination; and the hardware-specific Backend constructs a register interference graph to assign bounded physical registers via graph coloring.}
\label{ch41:fig-compiler-pipeline}
\end{figure}

% ======================================================================
% SECTION 1: LEXING
% ======================================================================
\section{Stage 1: Lexing --- From Character Streams to Tokens via Finite Automata}
\label{ch41:lexing}

The input to a compiler is a file on disk: a flat byte array of ASCII or UTF-8
characters. When your editor saves our running program, the file contains:
\[
  \texttt{s},\;\texttt{u},\;\texttt{m},\;\texttt{ },\;\texttt{=},\;\texttt{ },\;\texttt{0},\;\texttt{\textbackslash n},\;\texttt{w},\;\texttt{h},\;\texttt{i},\;\texttt{l},\;\texttt{e},\;\dots
\]
Why can the parser not read this character array directly?

If the parser operated on raw characters, every single grammar rule would have
to account for optional whitespace, tabs, line breaks, and comments between
every letter of every keyword and identifier. In our running example, the
grammar would need rules describing that \texttt{'s'} followed by \texttt{'u'}
followed by \texttt{'m'} forms the identifier \texttt{sum}, and would have to
permit arbitrary spaces between \texttt{s} and \texttt{u}. Such an approach
leads to a combinatorial explosion of grammatical rules.

The compiler solves this by delegating character scanning to Stage~1: the
\emph{lexical analyzer} or \emph{lexer}. The lexer scans the raw byte stream,
discarding whitespace and comments and grouping contiguous characters into
atomic units called \textbf{tokens}.

\begin{definition}[Token]
\label{ch41:def-token}
A \emph{token} is a pair $(\text{type}, \text{attribute})$, where $\text{type}$
is an enumerated category representing a syntactic role (such as an identifier,
keyword, or operator), and $\text{attribute}$ is the optional payload (such as
the string text of an identifier or the numerical value of a literal).
\end{definition}

Table~\ref{ch41:tab-tokens} displays the stream of tokens produced by the lexer
for our running program.

\begin{table}[H]
\centering
\small
\begin{tabular}{lll@{\qquad}lll}
\toprule
\textbf{Snippet} & \textbf{Token type} & \textbf{Attribute} & \textbf{Snippet} & \textbf{Token type} & \textbf{Attribute} \\
\midrule
\texttt{sum} & \texttt{IDENT} & \texttt{"sum"} & \texttt{+} & \texttt{OP\_PLUS} & --- \\
\texttt{=} & \texttt{OP\_ASSIGN} & --- & \texttt{a} & \texttt{IDENT} & \texttt{"a"} \\
\texttt{0} & \texttt{INT\_LIT} & $0$ & \texttt{[} & \texttt{LBRACKET} & --- \\
\texttt{while} & \texttt{KW\_WHILE} & --- & \texttt{i} & \texttt{IDENT} & \texttt{"i"} \\
\texttt{(} & \texttt{LPAREN} & --- & \texttt{]} & \texttt{RBRACKET} & --- \\
\texttt{i} & \texttt{IDENT} & \texttt{"i"} & \texttt{*} & \texttt{OP\_STAR} & --- \\
\texttt{<} & \texttt{OP\_LT} & --- & \texttt{2} & \texttt{INT\_LIT} & $2$ \\
\texttt{n} & \texttt{IDENT} & \texttt{"n"} & \texttt{i} & \texttt{IDENT} & \texttt{"i"} \\
\texttt{)} & \texttt{RPAREN} & --- & \texttt{=} & \texttt{OP\_ASSIGN} & --- \\
\texttt{\{} & \texttt{LBRACE} & --- & \texttt{i} & \texttt{IDENT} & \texttt{"i"} \\
\texttt{sum} & \texttt{IDENT} & \texttt{"sum"} & \texttt{+} & \texttt{OP\_PLUS} & --- \\
\texttt{=} & \texttt{OP\_ASSIGN} & --- & \texttt{1} & \texttt{INT\_LIT} & $1$ \\
\texttt{sum} & \texttt{IDENT} & \texttt{"sum"} & \texttt{\}} & \texttt{RBRACE} & --- \\
\bottomrule
\end{tabular}
\caption{The token stream produced by lexical analysis on our running example.}
\label{ch41:tab-tokens}
\end{table}

\subsection{The Maximal Munch Principle}

Tokenizing introduces boundary ambiguities. When encountering \texttt{<=}, does
the lexer emit \texttt{OP\_LE}, or \texttt{OP\_LT} followed by \texttt{OP\_ASSIGN}?
When reading \texttt{while1}, is it the keyword \texttt{while} followed by $1$,
or a single identifier?

Lexers resolve this with the \textbf{Maximal Munch} principle (the \emph{longest
match rule}):

\begin{formalbox}
\textbf{Maximal Munch Rule:} The lexer consumes the longest prefix of remaining
input that matches a valid token specification. When the longest match fits
multiple patterns (for example, \texttt{while} matches both keyword and identifier
rules), reserved keyword rules take precedence over general identifier rules.
\end{formalbox}

Under Maximal Munch, \texttt{<=} is recognized as a single \texttt{OP\_LE} token
because two characters form a longer match than one. Similarly, \texttt{while1}
is treated as an identifier.

\subsection{Scanning via Deterministic Finite Automata}

To execute Maximal Munch in strictly linear time, compilers compile regular
expressions into a \textbf{Deterministic Finite Automaton (DFA)}
(Chapter~31). A DFA consists of a state set $Q$, an alphabet
$\Sigma$, a transition function $\delta : Q \times \Sigma \to Q$, a start state
$q_0$, and accepting states $F \subseteq Q$.

\begin{algorithm}[H]
\caption{\textsc{DFA-Lexer-Next-Token}}
\KwIn{Character stream $S$, start index $p$, transition table $\delta$, accepting map $F$.}
\KwOut{Next token $T$, updated position $p$.}
$q \leftarrow q_0$; \quad $\text{last\_accepting\_state} \leftarrow \bot$; \quad $\text{last\_pos} \leftarrow p$; \quad $\text{scan\_pos} \leftarrow p$\;
\While{$\text{scan\_pos} \le |S|$ \textbf{and} $\delta(q, S[\text{scan\_pos}]) \neq \bot$}{
  $q \leftarrow \delta(q, S[\text{scan\_pos}])$\;
  $\text{scan\_pos} \leftarrow \text{scan\_pos} + 1$\;
  \If{$q \in F$}{
    $\text{last\_accepting\_state} \leftarrow q$; \quad $\text{last\_pos} \leftarrow \text{scan\_pos}$\;
  }
}
\If{$\text{last\_accepting\_state} = \bot$}{
  \Return $\text{Error: Unrecognized character at position } p$\;
}
$T \leftarrow \text{ConstructToken}(\text{last\_accepting\_state}, S[p\,..\,\text{last\_pos}-1])$\;
\Return $(T, \text{last\_pos})$\;
\end{algorithm}

Because the transition table $\delta[q, c]$ is precomputed in an array, each
character triggers an $O(1)$ state lookup. The entire source file of $n$ bytes
is tokenized in $O(n)$ time with zero backtracking.

% ======================================================================
% SECTION 2: PARSING
% ======================================================================
\section{Stage 2: Parsing --- From Flat Tokens to an Abstract Syntax Tree}
\label{ch41:parsing}

The lexer gives us a flat sequence of tokens. But computer programs are
intrinsically hierarchical: loops enclose blocks, blocks contain statements,
and statements enclose nested expressions.

\subsection{Why Finite Automata Cannot Parse Code}

Why not use finite automata to parse programming languages?

The limitation is mathematical. A finite automaton has a bounded number of
states: $|Q| = m$. Consider matching nested parentheses:
\[
  \underbrace{\texttt{((((}\dots\texttt{(}}_{k \text{ times}} a + b \underbrace{\texttt{)}\dots\texttt{)))}}
  _{k \text{ times}}
\]
To verify that the count of closing parentheses matches the opening ones, the
machine must store the current nesting depth. If the nesting depth $k$ exceeds
the state count $m$, the pigeonhole principle forces the automaton to enter some
state $q$ at least twice during the opening sequence. Having entered a cycle, the
automaton cannot distinguish $k$ open parentheses from $k + p$ open parentheses.
It cannot verify that the closing parentheses match.

Parsing nested syntax requires unbounded memory: a pushdown automaton with a
LIFO stack. The mathematical formalism for pushdown automata is the
\textbf{Context-Free Grammar (CFG)}.

\begin{definition}[Context-Free Grammar]
\label{ch41:def-cfg}
A \emph{Context-Free Grammar} is a 4-tuple $G = (V, \Sigma, R, S)$, where $V$ is a
finite set of nonterminal symbols, $\Sigma$ is the set of terminal tokens, $R$ is
a set of production rules of the form $A \to \alpha$ with $A \in V$ and
$\alpha \in (V \cup \Sigma)^*$, and $S \in V$ is the start symbol.
\end{definition}

\subsection{The Trap of Grammar Ambiguity}

Consider a naive expression grammar:
\[
  E \;\to\; E + E \;\;\mid\;\; E * E \;\;\mid\;\; \texttt{id} \;\;\mid\;\; \texttt{num}
\]
Now consider line~3 of our running program: $\texttt{sum} + \texttt{a[i]} * 2$.
Treating \texttt{a[i]} as $x$, this grammar allows two distinct parse trees:
\begin{enumerate}
  \item \textbf{Tree 1 (Addition first):} Apply $E \to E * E$ first, producing
  $(E + E) * E$, which computes $(\texttt{sum} + x) * 2$.
  \item \textbf{Tree 2 (Multiplication first):} Apply $E \to E + E$ first,
  producing $E + (E * E)$, which computes $\texttt{sum} + (x * 2)$.
\end{enumerate}

\begin{trapbox}
A grammar allowing two or more distinct parse trees for the same input string is
\textbf{ambiguous}. For $3 + 4 * 2$, Tree~1 evaluates to $(3 + 4) * 2 = 14$,
while Tree~2 evaluates to $3 + (4 * 2) = 11$. A compiler cannot leave operator
precedence to chance.
\end{trapbox}

\subsection{Resolving Ambiguity: Recursive Descent and Operator Precedence Climbing}

Compilers eliminate ambiguity by stratifying grammar rules into explicit
precedence tiers:
\begin{align*}
  E &\;\to\; E + T \;\;\mid\;\; T && \text{(Expressions: lowest precedence, addition)} \\
  T &\;\to\; T * F \;\;\mid\;\; F && \text{(Terms: medium precedence, multiplication)} \\
  F &\;\to\; \texttt{id} \;\;\mid\;\; \texttt{num} \;\;\mid\;\; \texttt{id}[E] \;\;\mid\;\; (E) && \text{(Factors: highest precedence, primary atoms)}
\end{align*}
In this stratified grammar, Tree~2 is the unique valid derivation: any
multiplication is parsed inside $T$, forcing $T * F$ to become a child of $E + T$.

In modern production compilers, this hierarchy is implemented via
\textbf{Recursive Descent} or Vaughan Pratt's 1973 \textbf{Operator Precedence
Climbing}. In Pratt parsing, every operator token has a numerical \emph{binding
power}. When parsing \texttt{sum + a[i] * 2}, the parser reads \texttt{+} (binding
power $10$). When the recursive sub-call encounters \texttt{*} (binding power
$20$), the higher binding power ($20 > 10$) pulls \texttt{a[i]} and \texttt{2}
into the multiplication subtree before the addition operator can close.

\subsection{The Abstract Syntax Tree (AST)}

Concrete parse trees retain punctuation tokens such as parentheses and semicolons.
Once grouping has been established, the compiler discards punctuation and
constructs an \textbf{Abstract Syntax Tree (AST)}. Interior nodes represent
operations (\texttt{While}, \texttt{Assign}, \texttt{BinOp}), and leaves
represent operands. Figure~\ref{ch41:fig-compiler-ast} illustrates the pipeline
from raw characters to tokens, and finally to an AST.

\begin{figure}[H]
\centering
\begin{tikzpicture}[
  font=\small,
  bytebox/.style={draw=gray!70, fill=gray!8, minimum width=5.5mm, minimum height=6.5mm, inner sep=0pt, font=\ttfamily\scriptsize},
  tokenpill/.style={draw=sourceblue!80, fill=sourceblue!10, rounded corners=3pt, inner sep=3pt, font=\ttfamily\scriptsize},
  astnode/.style={draw=structurepurple!80, fill=structurepurple!10, rounded corners=3pt, inner sep=3.5pt, font=\ttfamily\scriptsize, align=center},
  astleaf/.style={draw=targetgreen!80, fill=targetgreen!12, rounded corners=2pt, inner sep=3pt, font=\ttfamily\scriptsize},
  arrow/.style={-Stealth, thick, gray!60}, treeedge/.style={thick, gray!60}
]

% Panel (a): Raw ASCII Bytes
\node[anchor=west, font=\sffamily\bfseries\footnotesize, text=sourceblue!80!black] at (-5.8, 4.2) {(a) Raw ASCII Character Stream};
\foreach \c [count=\i from 0] in {s,u,m,{ },=,{ },0,$\backslash$n,w,h,i,l,e,{ },{(},i,{ },<,{ },n,{)},{ },{\{}} {
  \node[bytebox] at (-5.5 + \i*0.49, 3.5) {\c};
}
\node[anchor=west, font=\sffamily\scriptsize, text=gray!70] at (-5.5 + 23*0.49 + 0.1, 3.5) {$\dots$};

\draw[-Stealth, thick, gray!60] (0, 3.0) -- (0, 2.5) node[midway, right=3pt, font=\sffamily\scriptsize, text=gray!80] {Lexer (Finite Automaton)};

% Panel (b): Token Stream
\node[anchor=west, font=\sffamily\bfseries\footnotesize, text=sourceblue!80!black] at (-5.8, 2.1) {(b) Linear Token Stream};
\node[tokenpill] (tk1) at (-4.8, 1.4) {IDENT("sum")};
\node[tokenpill] (tk2) at (-2.8, 1.4) {OP\_ASSIGN};
\node[tokenpill] (tk3) at (-1.3, 1.4) {INT(0)};
\node[tokenpill] (tk4) at (0.2, 1.4) {KW\_WHILE};
\node[tokenpill] (tk5) at (1.6, 1.4) {LPAREN};
\node[tokenpill] (tk6) at (2.8, 1.4) {IDENT("i")};
\node[tokenpill] (tk7) at (4.0, 1.4) {OP\_LT};
\node[tokenpill] (tk8) at (5.2, 1.4) {IDENT("n")};

\draw[-Stealth, thick, gray!60] (0, 0.9) -- (0, 0.4) node[midway, right=3pt, font=\sffamily\scriptsize, text=gray!80] {Parser (Recursive Descent / Precedence Climbing)};

% Panel (c): AST
\node[anchor=west, font=\sffamily\bfseries\footnotesize, text=sourceblue!80!black] at (-5.8, 0.0) {(c) Abstract Syntax Tree (AST)};
\node[astnode] (root) at (0, -0.7) {Block};

% First child: sum = 0
\node[astnode] (asg1) at (-3.4, -1.8) {Assign};
\node[astleaf] (var_sum1) at (-4.4, -2.8) {sum}; \node[astleaf] (val_0) at (-2.4, -2.8) {0};
\draw[treeedge] (root) -- (asg1); \draw[treeedge] (asg1) -- (var_sum1); \draw[treeedge] (asg1) -- (val_0);

% Second child: While
\node[astnode] (wh) at (2.2, -1.8) {While}; \draw[treeedge] (root) -- (wh);
\node[astnode] (cond) at (0.4, -2.9) {BinOp($<$)};
\node[astleaf] (ci) at (-0.3, -3.9) {i}; \node[astleaf] (cn) at (1.1, -3.9) {n};
\draw[treeedge] (wh) -- node[left, font=\sffamily\tiny, text=gray!80] {cond} (cond);
\draw[treeedge] (cond) -- (ci); \draw[treeedge] (cond) -- (cn);

% While body
\node[astnode] (body) at (4.0, -2.9) {Block};
\draw[treeedge] (wh) -- node[right, font=\sffamily\tiny, text=gray!80] {body} (body);

% Body stmt 1: sum = sum + a[i] * 2
\node[astnode] (asg2) at (2.8, -4.0) {Assign};
\node[astleaf] (s2_dst) at (1.9, -5.0) {sum}; \node[astnode] (plus) at (3.5, -5.0) {BinOp($+$)};
\draw[treeedge] (body) -- (asg2); \draw[treeedge] (asg2) -- (s2_dst); \draw[treeedge] (asg2) -- (plus);

\node[astleaf] (plus_lhs) at (2.5, -6.1) {sum}; \node[astnode] (star) at (4.4, -6.1) {BinOp($\times$)};
\draw[treeedge] (plus) -- (plus_lhs); \draw[treeedge] (plus) -- (star);

\node[astnode] (idx) at (3.7, -7.2) {Index};
\node[astleaf] (idx_a) at (3.1, -8.1) {a}; \node[astleaf] (idx_i) at (4.1, -8.1) {i};
\draw[treeedge] (star) -- (idx); \draw[treeedge] (idx) -- (idx_a); \draw[treeedge] (idx) -- (idx_i);

\node[astleaf] (scale_2) at (5.1, -7.2) {2}; \draw[treeedge] (star) -- (scale_2);

% Body stmt 2: i = i + 1
\node[astnode] (asg3) at (5.7, -4.0) {Assign};
\node[astleaf] (i_dst) at (5.0, -5.0) {i}; \node[astnode] (i_plus) at (6.4, -5.0) {BinOp($+$)};
\draw[treeedge] (body) -- (asg3); \draw[treeedge] (asg3) -- (i_dst); \draw[treeedge] (asg3) -- (i_plus);

\node[astleaf] (i_lhs) at (5.8, -6.1) {i}; \node[astleaf] (i_rhs) at (7.0, -6.1) {1};
\draw[treeedge] (i_plus) -- (i_lhs); \draw[treeedge] (i_plus) -- (i_rhs);

\end{tikzpicture}
\caption{Tracing our running example across the compiler frontend: (a) raw character stream, (b) token stream after lexing, and (c) the resulting Abstract Syntax Tree (AST). Notice how operator precedence pushes $\text{BinOp}(\times)$ strictly deeper than $\text{BinOp}(+)$, eliminating typographical ambiguity without parentheses.}
\label{ch41:fig-compiler-ast}
\end{figure}

% ======================================================================
% SECTION 3: SEMANTIC ANALYSIS
% ======================================================================
\section{Stage 3: Semantic Analysis --- Scoped Symbol Tables and Type Lowering}
\label{ch41:semantics}

The AST in Figure~\ref{ch41:fig-compiler-ast} captures structure, but structure
alone does not guarantee validity. Consider the syntactically valid statement:
\[
  \texttt{x = y + 2}
\]
The parser accepts this statement without complaint. Yet at compile time, we must
know: Has \texttt{y} been declared? What is its type? If \texttt{y} is an open
network socket or a file handle, adding $2$ to it is invalid.

Stage~3, called \textbf{Semantic Analysis}, walks the AST to enforce semantic
validity and resolve variable bindings.

\subsection{Scoped Symbol Tables as a Stack of Hash Tables}

To track variable bindings across nested lexical blocks, compilers maintain a
\textbf{Symbol Table} using a stack of the \textbf{Hash Tables} introduced in
Chapter~11:
\begin{itemize}
  \item \textbf{Entering a scope:} When the AST walker enters a block delimited
  by \texttt{\{ \dots \}}, it pushes a fresh, empty hash table onto the scope stack.
  Each table retains a pointer to its enclosing parent table.
  \item \textbf{Declaring a symbol:} Variable declarations insert the identifier
  and its metadata (type, bit-width, memory offset) into the top hash table.
  \item \textbf{Resolving a symbol:} When an identifier appears in an expression,
  the compiler queries the top hash table in $O(1)$ expected time. If absent, it
  traverses parent pointers up to the global table. Inner variables naturally
  \emph{shadow} outer variables of the same name.
  \item \textbf{Exiting a scope:} Leaving a block pops the top table from the stack,
  instantly de-allocating local variables from visibility.
\end{itemize}

Figure~\ref{ch41:fig-symbol-tables} visualizes this stack discipline for our running example: nested syntactic blocks map directly to stacked hash tables connected by parent pointers.

\begin{figure}[H]
\centering
\begin{tikzpicture}[
  font=\small,
  scopebox/.style={draw=gray!60, rounded corners=4pt, inner sep=4pt},
  tablebox/.style={draw=gray!70, fill=white, rounded corners=3pt, minimum width=4.8cm, inner sep=4pt, align=left, font=\ttfamily\scriptsize},
  parentedge/.style={-Stealth, thick, gray!70, line width=1pt},
  resolvetrace/.style={-Stealth, dashed, thick, trapred!80!black}
]

% Panel (a): Nested Lexical Scopes
\node[anchor=west, font=\sffamily\bfseries\footnotesize, text=sourceblue!80!black] at (-6.5, 3.6) {(a) Nested Lexical Scopes};

\node[scopebox, fill=gray!8, minimum width=4.8cm, minimum height=4.8cm] (sc0) at (-4.1, 0.8) {};
\node[anchor=north west, font=\sffamily\bfseries\tiny, text=gray!80!black] at (-6.4, 3.0) {Scope 0: Global};
\node[anchor=north west, font=\ttfamily\scriptsize, align=left] at (-6.2, 2.6) {int64 a[], n;};

\node[scopebox, fill=sourceblue!10, draw=sourceblue!70, minimum width=4.2cm, minimum height=3.5cm] (sc1) at (-4.1, 0.4) {};
\node[anchor=north west, font=\sffamily\bfseries\tiny, text=sourceblue!90!black] at (-6.1, 1.95) {Scope 1: Function};
\node[anchor=north west, font=\ttfamily\scriptsize, align=left] at (-5.9, 1.55) {int64 sum = 0, i = 0;};

\node[scopebox, fill=targetgreen!12, draw=targetgreen!70, minimum width=3.6cm, minimum height=2.2cm] (sc2) at (-4.1, -0.05) {};
\node[anchor=north west, font=\sffamily\bfseries\tiny, text=targetgreen!90!black] at (-5.8, 0.85) {Scope 2: While Body};
\node[anchor=north west, font=\ttfamily\scriptsize, align=left] at (-5.6, 0.5) {int64 t1 = a[i];\\sum = sum + t1;\\i = i + 1;};

% Panel (b): Stack of Hash Tables
\node[anchor=west, font=\sffamily\bfseries\footnotesize, text=sourceblue!80!black] at (0.6, 3.6) {(b) Scope Stack \& Parent Chaining};

% Scope 2 Table (Top of Stack)
\node[tablebox, fill=targetgreen!8, draw=targetgreen!80] (tbl2) at (3.3, 2.1) {
  \textbf{\sffamily Scope 2 (Inner Block):}\\[1pt]
  \begin{tabular}{lll}
  \texttt{t1} & $\to$ & \textsf{type: int64,\; offset: -24}
  \end{tabular}
};

% Scope 1 Table
\node[tablebox, fill=sourceblue!8, draw=sourceblue!80] (tbl1) at (3.3, 0.5) {
  \textbf{\sffamily Scope 1 (Function Scope):}\\[1pt]
  \begin{tabular}{lll}
  \texttt{sum} & $\to$ & \textsf{type: int64,\; offset: -8}\\
  \texttt{i}   & $\to$ & \textsf{type: int64,\; offset: -16}
  \end{tabular}
};

% Scope 0 Table (Base of Stack)
\node[tablebox, fill=gray!6, draw=gray!70] (tbl0) at (3.3, -1.3) {
  \textbf{\sffamily Scope 0 (Global Scope):}\\[1pt]
  \begin{tabular}{lll}
  \texttt{a} & $\to$ & \textsf{type: array[int64],\; sym: \_a}\\
  \texttt{n} & $\to$ & \textsf{type: int64,\; sym: \_n}
  \end{tabular}
};

% Parent pointers
\draw[parentedge] (tbl2.south) -- node[midway, right=1pt, font=\sffamily\tiny, text=gray!70] {parent} (tbl1.north);
\draw[parentedge] (tbl1.south) -- node[midway, right=1pt, font=\sffamily\tiny, text=gray!70] {parent} (tbl0.north);

% Resolution trace for `a`
\draw[resolvetrace] (0.7, 2.1) to[bend right=22] node[midway, left=3pt, font=\sffamily\tiny, align=right] {lookup \texttt{a}:\\Scope 2 (miss)\\Scope 1 (miss)} node[pos=0.96, above left=1pt, font=\sffamily\tiny\bfseries, text=targetgreen!80!black] {HIT!} (tbl0.west);

\end{tikzpicture}
\caption{Scoped symbol tables organized as a stack of hash tables with enclosing parent pointers. In panel (a), nested blocks establish a lexical scope hierarchy. In panel (b), entering a scope pushes a new table with a parent link; identifier lookups probe the local table first and traverse parent pointers upward until hitting the binding or reaching the global scope.}
\label{ch41:fig-symbol-tables}
\end{figure}

\subsection{Type Checking and Lowering}

Once identifiers are bound, the compiler annotates each AST node with its static
type. In our running example:
\[
  \texttt{sum} : \text{int64},\quad \texttt{i} : \text{int64},\quad \texttt{n} : \text{int64},\quad \texttt{a} : \text{array}[\text{int64}].
\]
The semantic analyzer verifies that $\texttt{i} < \texttt{n}$ compares two
64-bit integers and that array elements match the type of \texttt{sum}.

Next comes \textbf{Type Lowering}. Modern silicon has no instruction that takes
an abstract array reference and returns an element. In hardware, array \texttt{a}
is simply a memory address pointer $\text{base}_a$. Because each 64-bit integer
requires $8$ bytes of storage, element $i$ sits at byte address:
\[
  \text{address}(a[i]) = \text{base}_a + i \times 8.
\]
The semantic analyzer lowers $\text{Index}(a, i)$ into an explicit pointer dereference:
\[
  \text{Index}(a, i) \;\Longrightarrow\; \text{load}\bigl(\text{base}_a + i \times 8\bigr).
\]
By making array indexing arithmetic explicit, the compiler exposes the hidden
multiplication by $8$ and the memory fetch to subsequent optimization passes.

% ======================================================================
% SECTION 4: IR, CFG, AND SSA
% ======================================================================
\section{Stage 4: Intermediate Representation, Control-Flow Graphs, and SSA Form}
\label{ch41:ir-ssa}

Translating $M$ source languages to $N$ computer architectures directly requires
$M \times N$ compilers. Production compilers avoid this by compiling frontends
to a common \textbf{Intermediate Representation (IR)}. With a universal IR, we
need only $M$ frontends and $N$ backends ($M + N$ components), and all
optimizations run on the shared IR.

\subsection{Three-Address Code and the Control-Flow Graph}

Modern IRs express computation as \textbf{Three-Address Code (3AC)}, where
nested expressions are broken into flat instructions with at most one operator:
$x \leftarrow y \text{ op } z$. To model control flow, instructions are grouped
into \textbf{Basic Blocks} connected in a \textbf{Control-Flow Graph (CFG)}
(Chapter~22).

\begin{definition}[Basic Block and CFG]
\label{ch41:def-cfg-block}
A \emph{Basic Block} is a maximal sequence of consecutive instructions where
control enters exclusively at the first instruction and leaves exclusively at
the last instruction. A \emph{Control-Flow Graph} is a directed graph $G = (V, E)$
where vertices $V$ are basic blocks and directed edges $(B_1, B_2) \in E$
represent possible branch or fallthrough transfers.
\end{definition}

\subsection{The Mutable Variable Dilemma and Static Single Assignment (SSA)}

In high-level languages, variables are mutable memory cells. In our loop,
\texttt{sum} is updated on every iteration. This mutability obscures dataflow:
to determine what value \texttt{sum} holds at an instruction, an optimizer must
trace all incoming execution paths.

In 1989, Cytron et al.\ introduced the standard representation for modern
optimizers: \textbf{Static Single Assignment (SSA)} form.

\begin{formalbox}
\textbf{The SSA Property:} A program is in \emph{Static Single Assignment form}
if and only if every variable is assigned a value at most once in the program text.
\end{formalbox}

Instead of overwriting variables, each assignment creates a uniquely numbered,
immutable version: $\text{sum}_0, \text{sum}_1, \text{sum}_2$. Each variable
becomes an immutable value, and definition-to-use relationships form explicit
edges in a Directed Acyclic Graph (DAG).

\subsection{The \texorpdfstring{$\phi$}{Phi} (Phi) Function at Control-Flow Joins}

At the top of our \texttt{while} loop, \texttt{sum} has value $\text{sum}_0 = 0$
on iteration~0 (from the entry block), but has value $\text{sum}_2$ on iteration
$k \ge 1$ (from the loop body). Because an SSA variable cannot be assigned twice,
how does the loop header refer to \texttt{sum}?

SSA introduces a \textbf{$\phi$ (phi) function}:
\[
  \text{sum}_1 = \phi\bigl(B_{\text{entry}} : \text{sum}_0,\;\; B_{\text{body}} : \text{sum}_2\bigr).
\]
A $\phi$-node is not a runtime instruction; microprocessors spend zero clock
cycles executing it. It is a compile-time marker that instructs the optimizer:
\emph{``If execution arrived from $B_{\text{entry}}$, use $\text{sum}_0$; if
execution arrived from $B_{\text{body}}$, use $\text{sum}_2$.''}

\subsection{Four Classic Optimizations Simplified by SSA}

Because SSA turns variable updates into explicit def-use graph edges, program
optimizations become clean graph rewrites:
\begin{enumerate}
  \item \textbf{Constant Propagation and Folding:} Constant operands are evaluated
  at compile time. Constant multiplications such as $t_2 \times 2$ undergo
  \emph{strength reduction} into single-cycle bitwise shifts: $t_2 \ll 1$.
  \item \textbf{Dead Code Elimination (DCE):} In SSA, every variable tracks its
  list of uses. If an SSA variable has zero uses and its defining instruction has
  no side effects, the instruction is deleted in $O(1)$ operations.
  \item \textbf{Common Subexpression Elimination (CSE) / Global Value Numbering (GVN):}
  Expressions are hashed into a Chapter~11 Hash Table keyed
  by $(\text{operator}, \text{val}_1, \text{val}_2)$. Duplicate computations are
  replaced with the precomputed value ID.
  \item \textbf{Loop-Invariant Code Motion (LICM) and Strength Reduction:}
  In our loop, $i_1$ is an induction variable stepping by $1$. The address
  $\text{base}_a + i_1 \times 8$ therefore steps by exactly $1 \times 8 = 8$ bytes
  per iteration. The compiler replaces the loop multiplication with a pointer
  variable $\text{ptr}$ that increments by $8$:
  \[
    \text{ptr}_0 = \text{base}_a,\quad \text{ptr}_2 = \text{ptr}_1 + 8.
  \]
  It then hoists the termination limit $\text{limit}_0 = \text{base}_a + (n_0 \ll 3)$
  into the entry block. The loop condition becomes $\text{ptr}_1 < \text{limit}_0$,
  and $i$ is eliminated by DCE!
\end{enumerate}

Figure~\ref{ch41:fig-ssa-cfg} shows the CFG before and after these transformations.

\begin{figure}[H]
\centering
\begin{tikzpicture}[
  font=\small,
  cfgblock/.style={draw=gray!70, fill=gray!4, rounded corners=3pt, inner sep=5pt, font=\ttfamily\scriptsize, align=left},
  headerblock/.style={draw=sourceblue!80, fill=sourceblue!6, rounded corners=3pt, inner sep=5pt, font=\ttfamily\scriptsize, align=left},
  optblock/.style={draw=targetgreen!80, fill=targetgreen!6, rounded corners=3pt, inner sep=5pt, font=\ttfamily\scriptsize, align=left},
  edge/.style={-Stealth, thick, gray!70}, trueedge/.style={-Stealth, thick, targetgreen!80!black},
  falseedge/.style={-Stealth, thick, trapred!80!black}, backedge/.style={-Stealth, thick, quotientorange!90!black}
]

% Panel A: Initial SSA CFG
\node[anchor=west, font=\sffamily\bfseries\footnotesize, text=sourceblue!80!black] at (-4.2, 4.8) {(a) Initial SSA Control-Flow Graph};
\node[cfgblock] (A_entry) at (-1.5, 3.5) {
  \textbf{\sffamily Block Entry:}\\
  $\text{sum}_0 = 0$;\; $i_0 = 0$\\
  \textsf{jump Header}
};
\node[headerblock] (A_header) at (-1.5, 1.4) {
  \textbf{\sffamily Block Header (Join):}\\
  $\text{sum}_1 = \phi(\text{Entry}\!:\!\text{sum}_0,\;\text{Body}\!:\!\text{sum}_2)$\\
  $i_1 = \phi(\text{Entry}\!:\!i_0,\;\text{Body}\!:\!i_2)$\\
  $\text{cond}_0 = (i_1 < n_0)$\\
  \textsf{br } $\text{cond}_0$ \textsf{ ? Body : Exit}
};
\node[cfgblock] (A_body) at (-1.5, -1.3) {
  \textbf{\sffamily Block Body:}\\
  $t_1 = i_1 \times 8$;\; $t_2 = \text{load}(\text{base}_a + t_1)$\\
  $t_3 = t_2 \times 2$;\; $\text{sum}_2 = \text{sum}_1 + t_3$\\
  $i_2 = i_1 + 1$;\; \textsf{jump Header}
};
\node[cfgblock] (A_exit) at (2.5, 1.4) {
  \textbf{\sffamily Block Exit:}\\
  \textsf{return } $\text{sum}_1$
};
\draw[edge] (A_entry) -- (A_header);
\draw[trueedge] (A_header) -- node[midway, right=2pt, font=\sffamily\tiny, text=targetgreen!80!black] {true} (A_body);
\draw[falseedge] (A_header) -- node[midway, above=2pt, font=\sffamily\tiny, text=trapred!80!black] {false} (A_exit);
\draw[backedge] (A_body.west) -- ++(-0.6, 0) |- (A_header.west)
  node[pos=0.25, left=2pt, font=\sffamily\tiny, text=quotientorange!90!black] {back-edge};

% Panel B: Optimized CFG
\node[anchor=west, font=\sffamily\bfseries\footnotesize, text=sourceblue!80!black] at (-4.2, -3.0) {(b) Optimized CFG (Strength Reduction \& Folding)};
\node[optblock] (B_entry) at (-1.5, -4.3) {
  \textbf{\sffamily Block Entry (Hoisted):}\\
  $\text{sum}_0 = 0$;\; $\text{ptr}_0 = \text{base}_a$\\
  $\text{limit}_0 = \text{base}_a + (n_0 \ll 3)$\\
  \textsf{jump Header}
};
\node[optblock] (B_header) at (-1.5, -6.5) {
  \textbf{\sffamily Block Header:}\\
  $\text{sum}_1 = \phi(\text{Entry}\!:\!\text{sum}_0,\;\text{Body}\!:\!\text{sum}_2)$\\
  $\text{ptr}_1 = \phi(\text{Entry}\!:\!\text{ptr}_0,\;\text{Body}\!:\!\text{ptr}_2)$\\
  $\text{cond}_0 = (\text{ptr}_1 < \text{limit}_0)$\\
  \textsf{br } $\text{cond}_0$ \textsf{ ? Body : Exit}
};
\node[optblock] (B_body) at (-1.5, -9.4) {
  \textbf{\sffamily Block Body (Reduced):}\\
  $\text{val}_0 = \text{load}(\text{ptr}_1)$\\
  $\text{scaled}_0 = \text{val}_0 \ll 1$;\; $\text{sum}_2 = \text{sum}_1 + \text{scaled}_0$\\
  $\text{ptr}_2 = \text{ptr}_1 + 8$;\; \textsf{jump Header}
};
\node[optblock] (B_exit) at (2.5, -6.5) {
  \textbf{\sffamily Block Exit:}\\
  \textsf{return } $\text{sum}_1$
};
\draw[edge] (B_entry) -- (B_header);
\draw[trueedge] (B_header) -- node[midway, right=2pt, font=\sffamily\tiny, text=targetgreen!80!black] {true} (B_body);
\draw[falseedge] (B_header) -- node[midway, above=2pt, font=\sffamily\tiny, text=trapred!80!black] {false} (B_exit);
\draw[backedge] (B_body.west) -- ++(-0.6, 0) |- (B_header.west)
  node[pos=0.25, left=2pt, font=\sffamily\tiny, text=quotientorange!90!black] {back-edge};

\end{tikzpicture}
\caption{The Control-Flow Graph (CFG) in Static Single Assignment (SSA) form: (a) initial SSA form with $\phi$-nodes at the loop header join point, and (b) after hoisting $\text{limit}_0$, reducing $i \times 8$ to pointer increment $\text{ptr} + 8$, and folding $t_2 \times 2 \to \text{val}_0 \ll 1$. All multiplications have been removed from the inner loop.}
\label{ch41:fig-ssa-cfg}
\end{figure}

% ======================================================================
% SECTION 5: THE BACKEND
% ======================================================================
\section{Stage 5: The Backend --- Instruction Selection, Liveness, and Register Allocation}
\label{ch41:backend}

The optimizer leaves our loop with an infinite supply of virtual variables. But
microprocessors have a strictly bounded bank of physical registers (16 on
x86-64, 31 on ARM64). Register accesses cost zero extra clock cycles; spilling a
variable to stack RAM incurs memory latency.

The backend maps virtual variables to $K$ physical registers without allowing
concurrently live values to overwrite each other.

\subsection{Liveness Analysis: Working Backward Through Time}

To assign registers safely, the compiler computes when each variable is active:

\begin{definition}[Live Variable]
\label{ch41:def-liveness}
A variable $v$ is \emph{live} at program point $p$ if there exists an execution
path from $p$ that reads $v$ before any redefinition of $v$.
\end{definition}

Liveness is computed \textbf{backward} from loop exit to entry:
\[
  \text{LiveIn}[B] = \text{Use}[B] \cup \bigl(\text{LiveOut}[B] \setminus \text{Def}[B]\bigr).
\]
Because the conditional back-edge at the bottom of the loop jumps back to the top
(and the exit block returns $\text{sum}$), the three loop-carried values
$\{\text{sum}, \text{ptr}, \text{limit}\}$ remain live across the entire loop,
whereas temporary values live only between their definition and their single read.
Tracing the five instructions of the optimized loop body:
\begin{enumerate}
  \item \textbf{Insn 1:} $\text{val} = \text{load}(\text{ptr})$ --- Live in: $\{\text{sum}, \text{ptr}, \text{limit}\}$. Live out: $\{\text{sum}, \text{ptr}, \text{limit}, \text{val}\}$.
  \item \textbf{Insn 2:} $\text{scaled} = \text{val} \ll 1$ --- Reads $\text{val}$; defines $\text{scaled}$. Variable $\text{val}$ is never read again and \textbf{dies}! Live out: $\{\text{sum}, \text{ptr}, \text{limit}, \text{scaled}\}$.
  \item \textbf{Insn 3:} $\text{sum} = \text{sum} + \text{scaled}$ --- Reads $\text{scaled}$; variable $\text{scaled}$ \textbf{dies}! Live out: $\{\text{sum}, \text{ptr}, \text{limit}\}$.
  \item \textbf{Insn 4:} $\text{ptr} = \text{ptr} + 8$ --- Updates $\text{ptr}$. Live out: $\{\text{sum}, \text{ptr}, \text{limit}\}$.
  \item \textbf{Insn 5:} $\text{cmp ptr, limit; jb loop}$ --- Compares $\text{ptr}$ and $\text{limit}$; jumps to header.
\end{enumerate}
Here $\text{val}$ dies at instruction~2 and $\text{scaled}$ is born at
instruction~2. Their live ranges \textbf{never overlap in time}.

\subsection{The Register Interference Graph and Chaitin's Graph Coloring}

\begin{definition}[Register Interference Graph]
\label{ch41:def-rig}
The \emph{Register Interference Graph (RIG)} is an undirected graph $G = (V, E)$
where vertices $V$ are program variables, and an edge $(u, v) \in E$ indicates
that $u$ and $v$ are simultaneously live at some program point.
\end{definition}

In 1981, Gregory Chaitin proved that assigning $K$ physical registers to variables
without conflict is isomorphic to \textbf{$K$-Vertex Coloring}
(Chapter~33). Although graph coloring is NP-complete,
compilers use the \textbf{Chaitin--Briggs degree-peeling heuristic}:
\begin{enumerate}
  \item \textbf{Simplify:} Find a vertex $v$ with $\deg(v) < K$. Push $v$ onto a
  stack and remove it from the graph. Because $v$ has at most $K-1$ neighbors,
  even if all neighbors receive distinct colors, at least one color among $\{1,
  \dots, K\}$ is guaranteed to remain free for $v$!
  \item \textbf{Spill:} If all vertices have degree $\ge K$, select a spill candidate
  for stack storage and continue peeling.
  \item \textbf{Select:} Pop vertices from the stack and assign each vertex any
  color unused by its already-colored neighbors.
\end{enumerate}

Tracing on our loop with $K=4$ registers ($R_0, R_1, R_2, R_3$):
\begin{itemize}
  \item $\text{val}$ and $\text{scaled}$ have $\deg = 3 < 4$.
  \item There is \textbf{no edge} between $\text{val}$ and $\text{scaled}$ because their live intervals are disjoint.
  \item Simplify peels $\text{scaled}$, then $\text{val}$, then the rest.
  \item The Select phase colors: $\text{sum} \mapsto R_0$, $\text{ptr} \mapsto R_1$, $\text{limit} \mapsto R_2$, $\text{val} \mapsto R_3$, and $\text{scaled} \mapsto R_3$.
\end{itemize}
Variables $\text{val}$ and $\text{scaled}$ share physical register $R_3$ with
zero conflicts! Figure~\ref{ch41:fig-register-coloring} illustrates the live
ranges and colored graph.

\begin{figure}[H]
\centering
\begin{tikzpicture}[
  font=\small,
  insnlabel/.style={anchor=east, font=\ttfamily\scriptsize, text=gray!90},
  vbar/.style={line width=3.5pt, line cap=round},
  rignode/.style={circle, draw=gray!70, minimum size=11.5mm, inner sep=1pt, font=\sffamily\scriptsize, align=center, thick},
  edge/.style={thick, gray!60}
]

% Panel A: Live Ranges
\node[anchor=west, font=\sffamily\bfseries\footnotesize, text=sourceblue!80!black] at (-5.8, 3.2) {(a) Variable Live-Range Intervals};
\foreach \y/\name in {
  2.4/{1:\;\;val = load(ptr)}, 1.6/{2:\;\;scaled = val $\ll$ 1},
  0.8/{3:\;\;sum = sum + scaled}, 0.0/{4:\;\;ptr = ptr + 8},
 -0.8/{5:\;\;cmp ptr, limit; jb loop}
} {
  \node[insnlabel] at (-2.3, \y) {\name};
  \draw[dashed, gray!25] (-2.1, \y) -- (1.1, \y);
}

\node[font=\ttfamily\bfseries\scriptsize, text=sourceblue!90!black] at (-1.8, 2.9) {sum};
\node[font=\ttfamily\bfseries\scriptsize, text=targetgreen!90!black] at (-1.1, 2.9) {ptr};
\node[font=\ttfamily\bfseries\scriptsize, text=quotientorange!90!black] at (-0.4, 2.9) {limit};
\node[font=\ttfamily\bfseries\scriptsize, text=structurepurple!90!black] at (0.3, 2.9) {val};
\node[font=\ttfamily\bfseries\scriptsize, text=structurepurple!90!black] at (1.0, 2.9) {scaled};

\draw[vbar, sourceblue!70] (-1.8, 2.6) -- (-1.8, -1.0);
\draw[vbar, targetgreen!70] (-1.1, 2.6) -- (-1.1, -1.0);
\draw[vbar, quotientorange!70] (-0.4, 2.6) -- (-0.4, -1.0);
\draw[vbar, structurepurple!70] (0.3, 2.4) -- (0.3, 1.6);
\draw[vbar, structurepurple!70] (1.0, 1.6) -- (1.0, 0.8);

\draw[stealth-stealth, semithick, structurepurple!80] (1.5, 2.4) -- (1.5, 1.6) node[midway, right=2pt, font=\sffamily\tiny] {live range $val$};
\draw[stealth-stealth, semithick, structurepurple!80] (1.5, 1.6) -- (1.5, 0.8) node[midway, right=2pt, font=\sffamily\tiny] {live range $scaled$};
\node[anchor=west, font=\sffamily\scriptsize, text=structurepurple!90!black] at (1.4, 0.2) {Disjoint intervals!};

% Panel B: RIG
\node[anchor=west, font=\sffamily\bfseries\footnotesize, text=sourceblue!80!black] at (3.8, 3.2) {(b) Interference Graph ($K=4$ Colored)};
\node[rignode, fill=sourceblue!25] (n_sum) at (5.8, 2.3) {$\text{sum}$\\[-2pt]\scriptsize\bfseries $R_0$};
\node[rignode, fill=targetgreen!25] (n_ptr) at (4.3, 1.0) {$\text{ptr}$\\[-2pt]\scriptsize\bfseries $R_1$};
\node[rignode, fill=quotientorange!25] (n_limit) at (4.8, -0.6) {$\text{limit}$\\[-2pt]\scriptsize\bfseries $R_2$};
\node[rignode, fill=structurepurple!25] (n_scaled) at (6.8, -0.6) {$\text{scaled}$\\[-2pt]\scriptsize\bfseries $R_3$};
\node[rignode, fill=structurepurple!25] (n_val) at (7.3, 1.0) {$\text{val}$\\[-2pt]\scriptsize\bfseries $R_3$};

\draw[edge] (n_sum) -- (n_ptr); \draw[edge] (n_sum) -- (n_limit);
\draw[edge] (n_sum) -- (n_val); \draw[edge] (n_sum) -- (n_scaled);
\draw[edge] (n_limit) -- (n_ptr); \draw[edge] (n_limit) -- (n_val); \draw[edge] (n_limit) -- (n_scaled);
\draw[edge] (n_ptr) -- (n_val); \draw[edge] (n_ptr) -- (n_scaled);
\draw[dashed, trapred!80, thick] (n_val) to[bend left=20] node[midway, right=1pt, font=\sffamily\tiny\bfseries, text=trapred!90!black] {no edge!} (n_scaled);

\end{tikzpicture}
\caption{Register allocation via Graph Coloring: (a) Live-range intervals across the five instructions of the optimized loop. Because $\text{val}$ dies at instruction~2 and $\text{scaled}$ is born at instruction~2, their live ranges are strictly disjoint. (b) The Register Interference Graph. Because $\text{val}$ and $\text{scaled}$ do not interfere, they share physical register $R_3$. The graph is validly 4-colored using registers $R_0, R_1, R_2, R_3$ with zero memory spills.}
\label{ch41:fig-register-coloring}
\end{figure}

\subsection{The Emitted Machine Code}

Mapping registers $R_0 \to \texttt{\%rax}$, $R_1 \to \texttt{\%rsi}$, $R_2 \to \texttt{\%rdx}$,
and $R_3 \to \texttt{\%rcx}$ (after a standard \emph{loop-rotation} entry guard checks $n > 0$ so the conditional branch sits at the bottom), the compiler emits x86-64 machine assembly:
\begin{lstlisting}[language={[x86masm]Assembler}]
.Lloop:
    movq    (%rsi), %rcx      # rcx (val) = load 8 bytes from ptr (rsi)
    shlq    $1, %rcx          # rcx (scaled) = rcx << 1
    addq    %rcx, %rax        # rax (sum) += scaled
    addq    $8, %rsi          # rsi (ptr) += 8
    cmpq    %rdx, %rsi        # compare ptr against limit (rdx)
    jb      .Lloop            # if ptr < limit, branch to .Lloop
\end{lstlisting}

A naive compiler emits over $18$ stack memory accesses per iteration, stalling
the execution pipeline. The optimizing compiler delivers a six-instruction inner loop
(five fused micro-ops) with zero stack spills and pure register execution.

% ======================================================================
% SUMMARY & EXERCISES
% ======================================================================
\section{Summary: The Pipeline of Transformations}
\label{ch41:summary}

Table~\ref{ch41:tab-summary} maps each compilation stage to its data structures,
algorithms, and textbook chapter foundations.

\begin{table}[H]
\centering
\small
\setlength{\tabcolsep}{4.5pt}
\begin{tabular}{lllll}
\toprule
\textbf{Stage} & \textbf{Input} & \textbf{Output} & \textbf{Algorithmic Engine} & \textbf{Book Chapter} \\
\midrule
1. Lexing & ASCII bytes & Token stream & Deterministic Finite Automata & Ch.~31 \\
2. Parsing & Token stream & AST & Recursive Descent / Pratt Parsing & Ch.~10, 12 \\
3. Semantics & AST & Typed AST & Scoped Symbol Tables / Type Lowering & Ch.~11 \\
4. IR / SSA & Typed AST & SSA CFG & Def-Use DAGs / Value Numbering & Ch.~11, 22 \\
5. Backend & SSA CFG & Assembly & Liveness Analysis / Graph Coloring & Ch.~15, 33 \\
\bottomrule
\end{tabular}
\caption{The five stages of a modern optimizing compiler.}
\label{ch41:tab-summary}
\end{table}

\begin{connectionbox}
\textbf{Where to go next: language abstractions under the hood.}
What an optimizer can simplify depends on how much type and dispatch information the
source language preserves until compile time. When you write a generic function in C++
or Rust, \emph{monomorphization} stamps out a concrete copy for each type, letting the
SSA middle-end inline calls and reduce loops to register arithmetic just as in our
six-instruction example; when you call through a C++ \texttt{virtual} method, Rust
\texttt{dyn Trait} fat pointer, or Go interface, the compiler instead emits an indirect
jump through a \emph{vtable} (virtual method table) whose target cannot be resolved until
runtime. The companion volume \emph{Programming Languages From the Ground Up} examines how
vtables, closures, algebraic data types, borrow checking, garbage collectors, and
compiler-generated \texttt{async}/\texttt{await} state machines lower onto this exact
compiler pipeline.
\end{connectionbox}

Across parallel algorithms, machine learning, systems architecture, and compiler
design, every algorithm we studied shared a classical premise: bits are strictly
0 or 1, gates compute Boolean functions, and execution moves through discrete
physical states. In Chapter~39, we turn to a different physical
model: quantum computing, where qubits exist in superposition and quantum
interference enables algorithms like Grover's search and Shor's factoring.

\section*{Exercises}
\addcontentsline{toc}{section}{Exercises}

\begin{exercise}[Tracing an AST with Operator Precedence]
Consider the arithmetic expression:
\[
  x = a + b * c - d / e.
\]
Assuming standard operator precedence ($*, / > +, -$) and left-associativity
for operators of equal precedence, draw the unambiguous Abstract Syntax Tree
produced by a Pratt parser. Which operation forms the root of the expression tree?
\end{exercise}

\begin{exercise}[Constructing SSA Form for Branching Control Flow]
Consider the following program with an \texttt{if-else} branch:
\begin{lstlisting}[language=C]
x = 10;
if (condition) {
    x = x + 5;
    y = x * 2;
} else {
    x = x - 3;
    y = 0;
}
z = x + y;
\end{lstlisting}
\begin{enumerate}
  \item Draw the Control-Flow Graph showing the four basic blocks: \texttt{Entry},
  \texttt{Then}, \texttt{Else}, and \texttt{Join}.
  \item Convert the program into Static Single Assignment (SSA) form by introducing
  subscripted variable versions ($x_0, x_1, \dots$).
  \item Write down the exact $\phi$-functions required in the \texttt{Join} block
  for variables $x$ and $y$.
\end{enumerate}
\end{exercise}

\begin{exercise}[Register Allocation via Graph 3-Coloring]
A compiler optimizes a loop containing five temporary variables: $u, v, w, x, y$.
Liveness analysis reveals the following symmetric interference pairs:
\begin{align*}
  (u, v),\quad (u, w),\quad (v, w),\quad (v, x),\quad (w, x),\quad (w, y),\quad (x, y).
\end{align*}
Suppose the CPU has $K=3$ physical registers: $\{R_0, R_1, R_2\}$.
\begin{enumerate}
  \item Compute the degree of each vertex in the interference graph.
  \item Apply the Chaitin--Briggs \texttt{Simplify} heuristic step by step. Which
  vertex can be peeled first? Show the order in which vertices are pushed onto
  the stack.
  \item Pop the vertices in the \texttt{Select} phase and assign registers from
  $\{R_0, R_1, R_2\}$ without conflict. Can the graph be colored with zero spills?
\end{enumerate}
\end{exercise}

\begin{exercise}[Induction Variable Strength Reduction on 2D Arrays]
In row-major order, an array element $M[i][j]$ in an $N \times N$ matrix of 64-bit
integers sits at byte address:
\[
  \text{address}(M[i][j]) = \text{base}_M + (i \times N + j) \times 8.
\]
Consider an inner loop where $i$ is fixed and $j$ runs from $0$ to $N-1$:
\begin{lstlisting}[language=C]
j = 0;
while (j < N) {
    sum = sum + M[i][j];
    j = j + 1;
}
\end{lstlisting}
Explain how an optimizing compiler applies loop-invariant code motion and
induction variable strength reduction to replace the two multiplications
($i \times N$ and $(i \times N + j) \times 8$) with a single pointer increment
inside the loop. Write the optimized SSA pseudocode for the loop body.
\end{exercise}

\begin{exercise}[The Cost of Register Spilling]
Suppose a microprocessor takes $0$ extra cycles to read a physical register,
$4$ cycles to read from L1 data cache, and $200$ cycles on an L1 cache miss that
reads DRAM memory.
\begin{enumerate}
  \item If a loop body executes $10^7$ times and spills two variables to the
  stack, compute the minimum latency penalty (in clock cycles) incurred by the
  spill assuming all stack accesses hit the L1 cache.
  \item Why do compiler register allocators use loop nest depth as a primary
  weighting factor when selecting which variable to spill?
\end{enumerate}
\end{exercise}
